% GENERATED FILE. Edit canonical lessons, then run npm run generate:books.
% canonical-source-hash: fnv1a64:935a8576372d0ea0
% canonical-lessons: LA-C311-damno, LA-C311-absolvo, LA-C311-nocens, LA-C311-innocens, LA-C311-culpa

\chapter{Damnātus aut absolūtus}
\label{ch:la-damnatus-aut-absolutus}
\hlchaptermodality{311}

\begin{chapteropening}
\textbf{By the end of this chapter:} \emph{I can tell the end of a trial: to condemn, to acquit, guilty, innocent and the blame.}

Complete the last lesson of chapter 311: I can tell the end of a trial: to condemn, to acquit, guilty, innocent and the blame.
\end{chapteropening}

\section[damnō, damnāre]{damnō, damnāre — to condemn, to find guilty}

\label{lesson:LA-C311-damno}



\begin{quote}
Before the new one: say the Latin for to prove, then the Latin for information.
\end{quote}

\subsection*{You'll want to know: damnō, damnāre}
\textbf{damnō, damnāre} — \textquotedblleft{}to condemn, to find guilty\textquotedblright{}.

\begin{quote}
Iūdicēs eum fūrtī damnāvērunt.
\end{quote}

\begin{quote}
\emph{The judges found him guilty of theft.}
\end{quote}

\emph{damnō, damnāre, damnāvī, damnātum}, with the genitive of the charge. From \emph{damnum}, 'loss' (taught). English \emph{damn}.

\begin{cousinweb}
\textbf{damnō} is built on \emph{damnum}, \textquotedblleft{}loss, damage\textquotedblright{} (EtymDB): to inflict a loss, a penalty. English \emph{damn} and \emph{condemn} come from its family.
\end{cousinweb}

\begin{grammarlens}[title={What you've built}]
One more word for telling, explaining and describing at B1.
\end{grammarlens}

\subsection*{Your turn}
\begin{itemize}
\raggedright
  \item \emph{Say it:} \emph{damnō, damnāre}
  \item \emph{Say it:} \emph{damnō, damnāre}, once more
  \item \emph{Say it:} a sentence of your own with \emph{damnō, damnāre}
  \item \emph{Recall:} say the Latin for to prove, then the Latin for information, then say \emph{damnō, damnāre} again
\end{itemize}

\begin{tcolorbox}[breakable,colback=teal!4,colframe=teal!35!black,title={Before you move on}]
What does damnō, damnāre mean? (\textquotedblleft{}To condemn, to find guilty\textquotedblright{}.) Say it once more.
\end{tcolorbox}
\section[absolvō, absolvere]{absolvō, absolvere — to acquit; to set free; to finish}

\label{lesson:LA-C311-absolvo}



\begin{quote}
Before the new one: say the Latin for information, then the Latin for to condemn.
\end{quote}

\subsection*{You'll want to know: absolvō, absolvere}
\textbf{absolvō, absolvere} — \textquotedblleft{}to acquit; to set free; to finish\textquotedblright{}.

\begin{quote}
Reus absolūtus est et laetus domum rediit.
\end{quote}

\begin{quote}
\emph{The defendant was acquitted and went home happy.}
\end{quote}

\emph{absolvō, absolvere, absolvī, absolūtum}: \emph{ab} + \emph{solvō} (taught), to loosen from. English \emph{absolve}.

\begin{cousinweb}
\textbf{absolvō} joins \emph{ab}, \textquotedblleft{}from\textquotedblright{}, and \emph{solvō}, \textquotedblleft{}to loosen\textquotedblright{} (EtymDB): to set loose from a charge. English \emph{absolve} comes from it.
\end{cousinweb}

\begin{grammarlens}[title={What you've built}]
One more word for telling, explaining and describing at B1.
\end{grammarlens}

\subsection*{Your turn}
\begin{itemize}
\raggedright
  \item \emph{Say it:} \emph{absolvō, absolvere}
  \item \emph{Say it:} \emph{absolvō, absolvere}, once more
  \item \emph{Say it:} a sentence of your own with \emph{absolvō, absolvere}
  \item \emph{Recall:} say the Latin for information, then the Latin for to condemn, then say \emph{absolvō, absolvere} again
\end{itemize}

\begin{tcolorbox}[breakable,colback=teal!4,colframe=teal!35!black,title={Before you move on}]
What does absolvō, absolvere mean? (\textquotedblleft{}To acquit; to set free; to finish\textquotedblright{}.) Say it once more.
\end{tcolorbox}
\section[nocēns, nocentis]{nocēns, nocentis — guilty; harmful}

\label{lesson:LA-C311-nocens}



\begin{quote}
Before the new one: say the Latin for to condemn, then the Latin for to acquit.
\end{quote}

\subsection*{You'll want to know: nocēns, nocentis}
\textbf{nocēns, nocentis} — \textquotedblleft{}guilty; harmful\textquotedblright{}.

\begin{quote}
Nocēns poenam dare dēbet.
\end{quote}

\begin{quote}
\emph{The guilty man must pay the penalty.}
\end{quote}

One ending, genitive \emph{nocentis}: the present participle of a verb meaning 'to harm'.

\begin{cousinweb}
\textbf{nocēns} is the present participle of \emph{noceō}, \textquotedblleft{}to harm\textquotedblright{} (Lewis \& Short): doing harm, and so guilty.
\end{cousinweb}

\begin{grammarlens}[title={What you've built}]
One more word for telling, explaining and describing at B1.
\end{grammarlens}

\subsection*{Your turn}
\begin{itemize}
\raggedright
  \item \emph{Say it:} \emph{nocēns, nocentis}
  \item \emph{Say it:} \emph{nocēns, nocentis}, once more
  \item \emph{Say it:} a sentence of your own with \emph{nocēns, nocentis}
  \item \emph{Recall:} say the Latin for to condemn, then the Latin for to acquit, then say \emph{nocēns, nocentis} again
\end{itemize}

\begin{tcolorbox}[breakable,colback=teal!4,colframe=teal!35!black,title={Before you move on}]
What does nocēns, nocentis mean? (\textquotedblleft{}Guilty; harmful\textquotedblright{}.) Say it once more.
\end{tcolorbox}
\section[innocēns, innocentis]{innocēns, innocentis — innocent, harmless}

\label{lesson:LA-C311-innocens}



\begin{quote}
Before the new one: say the Latin for to acquit, then the Latin for guilty.
\end{quote}

\subsection*{You'll want to know: innocēns, innocentis}
\textbf{innocēns, innocentis} — \textquotedblleft{}innocent, harmless\textquotedblright{}.

\begin{quote}
Omnēs scīvērunt eum innocentem esse.
\end{quote}

\begin{quote}
\emph{Everyone knew that he was innocent.}
\end{quote}

One ending, genitive \emph{innocentis}: \emph{in-} 'not' + \emph{nocēns}. English \emph{innocent}.

\begin{cousinweb}
\textbf{innocēns} joins \emph{in-}, \textquotedblleft{}not\textquotedblright{}, and \emph{nocēns}, \textquotedblleft{}harmful, guilty\textquotedblright{} (EtymDB): doing no harm. English \emph{innocent} comes from it.
\end{cousinweb}

\begin{grammarlens}[title={What you've built}]
One more word for telling, explaining and describing at B1.
\end{grammarlens}

\subsection*{Your turn}
\begin{itemize}
\raggedright
  \item \emph{Say it:} \emph{innocēns, innocentis}
  \item \emph{Say it:} \emph{innocēns, innocentis}, once more
  \item \emph{Say it:} a sentence of your own with \emph{innocēns, innocentis}
  \item \emph{Recall:} say the Latin for to acquit, then the Latin for guilty, then say \emph{innocēns, innocentis} again
\end{itemize}

\begin{tcolorbox}[breakable,colback=teal!4,colframe=teal!35!black,title={Before you move on}]
What does innocēns, innocentis mean? (\textquotedblleft{}Innocent, harmless\textquotedblright{}.) Say it once more.
\end{tcolorbox}
\section[culpa, culpae]{culpa, culpae — fault, blame; guilt}

\label{lesson:LA-C311-culpa}



\begin{quote}
Before the new one: say the Latin for guilty, then the Latin for innocent.
\end{quote}

\subsection*{You'll want to know: culpa, culpae}
\textbf{culpa, culpae} — \textquotedblleft{}fault, blame; guilt\textquotedblright{}.

\begin{quote}
Nōn mea culpa erat, sed tua.
\end{quote}

\begin{quote}
\emph{It was not my fault, but yours.}
\end{quote}

Feminine. \emph{In culpā esse}: to be to blame. English \emph{culpable} comes from it.

\begin{cousinweb}
\textbf{culpa} is inherited from Proto-Italic \emph{*k\textsuperscript{w}olpā}, \textquotedblleft{}a wrong, a mistake\textquotedblright{} (EtymDB). English \emph{culpable} and \emph{culprit} come from its family.
\end{cousinweb}

\begin{grammarlens}[title={What you've built}]
One more word for telling, explaining and describing at B1.
\end{grammarlens}

\subsection*{Your turn}
\begin{itemize}
\raggedright
  \item \emph{Say it:} \emph{culpa, culpae}
  \item \emph{Say it:} \emph{culpa, culpae}, once more
  \item \emph{Say it:} a sentence of your own with \emph{culpa, culpae}
  \item \emph{Recall:} say the Latin for guilty, then the Latin for innocent, then say \emph{culpa, culpae} again
\end{itemize}

\begin{tcolorbox}[breakable,colback=teal!4,colframe=teal!35!black,title={Before you move on}]
What does culpa, culpae mean? (\textquotedblleft{}Fault, blame; guilt\textquotedblright{}.) Say it once more.
\end{tcolorbox}
